% ============================================================================
\section{Discussion}\label{sec:discussion}
% ============================================================================

\subsection{Where consent is lost}\label{sec:disc-where}

Shopify's gate is not the weak point: it held back almost every configured pixel whose declared purpose had been refused, with exceptions on 7 of 1,179 stores. The interfaces around it are (Section~\ref{sec:bg-threat}). \emph{F1:} before any choice, the region default at our California vantage point allowed marketing on most stores, and sandboxed pixels fired on 70.2\% of banner stores with that default. \emph{F2:} where a blocklist declared an ad vendor blocked and the CMP had not itself consented, that vendor's pixel fired from the sandbox on 46 of 89 stores, and blocklist stores fired more often than comparable stores, because the sandbox followed Shopify's state, not the blocklist. \emph{F3:} marketing was still allowed after 20 of 231 working Decline clicks, automatic ``no'' writes arrived after 525 of 535 third-party marketing sandboxes had started, pixels started before a recorded refusal kept sending on 32 of 196 stores, and sale of data stayed allowed after 112 of the 196 refusals. \emph{F4:} pixels that declare no purpose kept sending after a recorded refusal on 28 of the 196 stores and under GPC on 191 of 1,354 and under a denying default they fired before any choice on 19 of 167 banner stores.

In security terms, the page carries two policies for the same code: the CMP's configuration states one (``TikTok: blocked'', for main-thread tags), and the platform enforces another (``allowed until the visitor declines, for the purposes the pixel declares''). The design departs from complete mediation, that every access be checked at the point that controls it \citep{saltzer1975protection}: the CMP checks tags on the main thread, but the WPM decides whether the code runs, and the WPM's own check rests on labels chosen by the party it restrains (Section~\ref{sec:bg-threat}), so a pixel that declares nothing is never checked. This gap between stated and enforced policy, of a kind security research examines elsewhere on the web \citep{latta2026policy}, follows from the architecture, not from one vendor, and may grow as more tracking moves into the sandbox \citep{shopify_checkout_upgrade}.

\subsection{Who can fix it}\label{sec:disc-fix}

\subsubsection*{Shopify}
The platform controls the only enforcement point that sees sandboxed pixels. (i) When a third-party banner is present, Shopify could hold the gate closed until the CMP has reported an initial state, as Google Consent Mode's \code{wait\_for\_update} setting does \citep{google_consentmode}, removing the race of Section~\ref{sec:res-race}. (ii) On a withdrawal of consent, the WPM could terminate the worker or remove the iframe of every pixel whose declared purpose is no longer allowed, or at least document and verifiably stop delivering events to it. (iii) Advertising apps could be required to declare the purposes their data flows serve, and merchants to declare those of custom pixels, which were involved in most residuals after a refusal; under GPC, Shopify could also hold back pixels that declare marketing, since the CCPA's opt-out covers ``sharing'' for cross-context behavioural advertising. (iv) Consent strings generated for one region should not be served to visitors from another; caching providers and merchants must keep the consent header out of cached responses.

\subsubsection*{CMP vendors}
Several vendors already document that their blockers cannot reach app or custom pixels (Section~\ref{sec:related}). They could (i) write the initial state and every choice to Shopify's Customer Privacy API by default, before the WPM's first decision; (ii) record a US Decline as a refusal of sale of data as well as marketing; (iii) show sandboxed pixels in their scanners as governed by Shopify's state, not by the blocker; (iv) offer merchants a self-test that checks, after a Decline, that no new sandbox sends ad requests; and (v) not grant consent to automated browsers, which misleads compliance scanners too.

\subsubsection*{Merchants}
A merchant who shows US visitors a banner offering a choice can make the platform consistent with it by setting the region to require consent, as almost all stores with Shopify's own banner had done, and by declaring the purposes that custom pixels serve.

\subsubsection*{Browsers and standards}
For persistence after the choice, GPC's advantage over a working Decline was at most about 10 percentage points, but GPC needs no banner integration, applies before any pixel starts, and will reach far more visitors under the 2027 browser requirement \citep{ab566,gpc_site}; it will not reach pixels that do not declare sale of data. The GPC specification exposes the signal to workers \citep{w3c_gpc}, but Shopify's STRICT sandbox hides it; preserving \code{globalPrivacyControl} inside such sandboxes would let vendor code honour it directly.

\subsection{A self-test for merchants}\label{sec:disc-tool}

Merchants cannot see these failures from their CMP's dashboard, which lists the blocked vendors, or from Shopify's, which lists the installed pixels. We therefore packaged the measurement as ConsentCheck, a command-line tool that a merchant runs against their own store: it reuses the crawler's attribution code, makes three visits with fresh profiles (Default, Decline followed by one reload, and GPC), and reports per vendor what left before any choice, from where, what the choices stopped and which blocklisted vendors the sandbox still runs, followed by the fixes of Section~\ref{sec:disc-fix} (safeguards and validation: Supplementary \ref*{S:sec:s-extended}). It is released with the code at \url{https://github.com/devanshsanghavi-droid/FTJ}.

\subsection{Legal and policy context}\label{sec:disc-law}

We make no claim that any store, CMP or platform acts unlawfully, and the measurements do not establish it; pre-choice firing is not by itself prohibited by the CCPA's opt-out rules (Section~\ref{sec:bg-law}), and we did not assess other obligations. Opt-out tools were not fully effective, mostly because of pixels whose declarations the visitor cannot see; whether those flows are a ``sale'' or ``sharing'' depends on contracts and data uses that a browser cannot observe. California enforcement has already involved opt-outs that tracking did not follow \citep{cppa2025toddsnyder,caag2025healthline,cppa2025tractorsupply}, and tracking that starts before consent has been litigated under the California Invasion of Privacy Act, including against Shopify \citep{briskin2025shopify,briskin2026remand}; SB~690, signed on 30 September 2026 (Chapter 976, Statutes of 2026) and operative from 1 January 2027, leaves suits against private actors under that act's pen-register provision over website and app tracking to the Attorney General \citep{leginfo2026sb690,duanemorris2026sb690}, and \citet{munir2025keystroke} map web tracking techniques to wiretap law (further detail: Supplementary \ref*{S:sec:s-extended}). EU visitors must consent before tracking \citep{eprivacy2002,edpb2023guidelines}; Shopify's EEA defaults may address F1 there, but our opt-in proxy (Section~\ref{sec:res-prechoice}) suggests that F4 would remain.

\subsection{Beyond Shopify}\label{sec:disc-beyond}

The blind spot is a property of any architecture that moves third-party code out of the main document: worker-offloading libraries \citep{partytown}, tag managers with their own consent layers \citep{mertens2025gtm} and server-side tagging \citep{bottger2026fpt,jazlan2026sstguard} all relocate the decision to run tracking code, and consent must then be enforced by the component that starts the code. Our pilot (Section~\ref{sec:res-beyond}) shows the consequence on other platforms: ad requests continued after a working Decline on 42\% of the Wix and 60\% of the BigCommerce sites we loaded, and first-party gateways kept collecting after refusals, although the pilot cannot say how often this happens across either platform or what a gateway forwards. For measurement the lesson is concrete: network-level crawlers see sandboxed requests but cannot tie them to the gate that started them without attribution, and audits that infer tracking from cookies \citep{bollinger2022automating,rasaii2025intractable,libert2026webxray} miss them: no ad request from a STRICT worker in the Default load carried a cookie (Supplementary \ref*{S:sec:s-extended}). Our join of declared purpose, platform consent state and container start can audit any consent-gated tag system, and we release it with our code. Automated consent audits should also report whether consent tools treat their browser as a bot.
